\section{补充总图：从抽象层回到完整栈}

前面已经把课程为什么存在、为什么强调效率、为什么要把规模看成一等公民讲清楚了。这里再把这门课的“完整栈”重新压成一张图，方便后续章节反复引用。

\begin{figure}[H]
\centering
\begin{tikzpicture}[node distance=1.0cm, every node/.style={font=\small, align=center}]
\node[draw, rounded corners, fill=blue!10, minimum width=3.8cm, minimum height=0.8cm] (a) {abstraction / API};
\node[draw, rounded corners, fill=blue!16, below=of a, minimum width=3.8cm, minimum height=0.8cm] (b) {leaky abstractions};
\node[draw, rounded corners, fill=blue!22, below=of b, minimum width=3.8cm, minimum height=0.8cm] (c) {mechanics};
\node[draw, rounded corners, fill=blue!28, below=of c, minimum width=3.8cm, minimum height=0.8cm] (d) {systems};
\node[draw, rounded corners, fill=blue!34, below=of d, minimum width=3.8cm, minimum height=0.8cm] (e) {data / scale / alignment};
\draw[-{Latex[length=3mm]}, thick] (a) -- (b);
\draw[-{Latex[length=3mm]}, thick] (b) -- (c);
\draw[-{Latex[length=3mm]}, thick] (c) -- (d);
\draw[-{Latex[length=3mm]}, thick] (d) -- (e);
\end{tikzpicture}
\caption{课程反复要你做的事情：从表层接口回到机制，再回到可改进的完整栈}
\end{figure}

\begin{importantbox}{完整栈视角}
这门课不是分别讲“模型”“系统”“数据”“对齐”，而是把它们看成一个耦合系统：哪怕你只改一个环节，其他环节的最优解也会变。
\end{importantbox}

\begin{figure}[H]
\centering
\begin{tikzpicture}[node distance=0.85cm, every node/.style={font=\small, align=center}]
\node[draw, rounded corners, fill=green!10, minimum width=3.2cm, minimum height=0.8cm] (build) {build};
\node[draw, rounded corners, fill=green!18, below=of build, minimum width=3.2cm, minimum height=0.8cm] (measure) {measure};
\node[draw, rounded corners, fill=green!26, below=of measure, minimum width=3.2cm, minimum height=0.8cm] (scale) {scale};
\node[draw, rounded corners, fill=green!34, below=of scale, minimum width=3.2cm, minimum height=0.8cm] (align) {align};
\draw[-{Latex[length=3mm]}, thick] (build) -- (measure);
\draw[-{Latex[length=3mm]}, thick] (measure) -- (scale);
\draw[-{Latex[length=3mm]}, thick] (scale) -- (align);
\end{tikzpicture}
\caption{课程的四步闭环：先构建，再测量，再放大，最后对齐到人的目标}
\end{figure}

\begin{knowledgebox}{本讲结束时应该形成的直觉}
\begin{itemize}
    \item 模型不是孤立对象，而是资源受限系统的一部分。
    \item 设计不是一次性的，而是 build-measure-scale-align 的循环。
    \item 研究不仅要问“能不能做”，还要问“值不值得做”“能不能继续做大”。
\end{itemize}
\end{knowledgebox}

